\subsection{根的闭集}

定义 4. 设 P 是 R 的子集。

\begin{enumerate}
\def\labelenumi{(\roman{enumi})}
\item
  若 \(\alpha \in \mathrm{P},\beta \in \mathrm{P},\alpha +\beta \in \mathrm{R}\)，则 \(\alpha +\beta \in \mathrm{P}\)，称 P 为闭的。
\item
  若 P 是闭的且 \(\mathrm{P} \cup (-\mathrm{P}) = \mathrm{R}\)，称 P 为抛物的。
\item
  若 \(\mathrm{P} = -\mathrm{P}\)，称 P 为对称的。
\end{enumerate}

引理 3. 设 C 是 R 的一个室，P 是包含 \(\mathrm{R}_{+}(\mathrm{C})\) 的 R 的闭子集（采用第 6 号的记号）。设 \(\Sigma = \mathrm{B}(\mathrm{C}) \cap (-\mathrm{P})\)，设 Q 为由 \(\Sigma\) 中元素以非正整数为系数组合成的根的集合。于是，\(\mathrm{P} = \mathrm{R}_{+}(\mathrm{C}) \cup \mathrm{Q}\)。

只需证明 \(P \cap (-R_{+}(C)) = Q\)。设 \(-\alpha \in Q\)。于是 \(\alpha\) 是 \(\Sigma\) 中 n 个元素之和。我们对 n 归纳证明 \(-\alpha \in P\)。当 n = 1 时显然成立。若 n \textgreater{} 1，则根据第 6 号的命题 19，可以写成 \(\alpha = \beta + \gamma\)，其中 \(\gamma \in \Sigma\)，而 \(\beta\) 是 \(\Sigma\) 中 n - 1 个元素之和。由归纳假设，\(-\beta \in P\)；由于 \(-\gamma \in P\) 且 P 是闭的，\(-\alpha \in P\)。因此，\(Q \subseteq P \cap (-R_{+}(C))\)。反过来，设 \(-\alpha \in P \cap (-R_{+}(C))\)。于是 \(\alpha\) 是 B(C) 中 p 个元素之和。我们对 p 归纳证明 \(-\alpha \in Q\)。当 p = 1 时显然成立。若 p \textgreater{} 1，则根据命题 19，可以写成 \(\alpha = \beta + \gamma\)，其中 \(\gamma \in B(C)\)，而 \(\beta\) 是 B(C) 中 p - 1 个元素之和。由于 \(-\gamma = \beta + (-\alpha)\) 且 P 是闭的，\(-\gamma \in P\)，从而 \(\gamma \in \Sigma\)。此外，\(-\beta = \gamma + (-\alpha) \text{ 故 } -\beta \in \mathbf{P} \text{ 因为 P 是闭的。由归纳假设， } -\beta \in \mathbf{Q}, \text{ 故 } -\alpha = -\beta - \gamma \in \mathbf{Q}\)。因此，\(\mathrm{P} \cap (-\mathrm{R}_{+}(\mathrm{C})) \subseteq \mathrm{Q}\)。

命题 20. 设 P 是 R 的子集。以下条件等价：

\begin{enumerate}
\def\labelenumi{(\roman{enumi})}
\item
  P 是抛物的；
\item
  P 是闭的，并且存在 R 的一个室 C，使得 \(\mathrm{P} \supseteq \mathrm{R}_{+}(\mathrm{C})\)；
\item
  存在 R 的一个室 C 以及 B(C) 的子集 \(\Sigma\)，使得 P 是 \(\mathrm{R}_{+}(\mathrm{C})\) 与根集 Q 的并集，其中 Q 是由 \(\Sigma\) 中元素以非正整数为系数组合成的根的集合。
\item
  \(\Longrightarrow\) (iii)：这由引理 3 得到。
\item
  \(\Longrightarrow\) (i)：采用 (iii) 的假设和记号。显然 \(\mathrm{P} \cup (-\mathrm{P}) = \mathrm{R}\)。我们证明，若 \(\alpha, \beta \in \mathrm{P}\) 且 \(\alpha + \beta \in \mathrm{R}\)，则 \(\alpha + \beta \in \mathrm{P}\)。若根 \(\alpha + \beta\) 为正，这显然成立。假设 \(\alpha + \beta\) 为负。于是 \(\alpha + \beta = \sum_{\gamma \in \mathrm{B}(\mathbb{C})} n_{\gamma} \gamma\)，其中 \(n_{\gamma} \leqslant 0\)。但 \(\gamma\) 是 \(\mathrm{B}(\mathbb{C}) - \Sigma\) 中的任意元素；它在 \(\alpha\) 或 \(\beta\) 中的系数都为 \(\geqslant 0\)；因此 \(n_{\gamma} = 0\) 当 \(\gamma \in \mathrm{B}(\mathbb{C}) - \Sigma\) 时，所以 \(\alpha + \beta \in \mathbb{Q} \subseteq \mathbb{P}\)。
\item
  \(\Longrightarrow\) (ii)：假设 P 是抛物的。设 C 是一个室，使 \(\operatorname{Card}(P\cap R_{+}(C))\) 尽可能大。令 \(\alpha\in\operatorname{B}(C)\)，并假设 \(\alpha\in P\)，于是 \(-\alpha\in P\)。对于所有 \(\beta\in P\cap R_{+}(C)\)，\(\beta\) 不与 \(\alpha\) 成比例（因为假设 \(\beta=2\alpha\) 会蕴含 \(\alpha=2\alpha+(-\alpha)\in P\)，这是由于 P 是闭的）。因此 \(s_{\alpha}(\beta)\in\operatorname{R}_{+}(C)\)（第 6 号的 Cor. 1 of Prop. 17）。若置 \(C'=s_{\alpha}(C)\)，则 \(\beta=s_{\alpha}(s_{\alpha}(\beta))\in s_{\alpha}(\operatorname{R}_{+}(C))=\operatorname{R}_{+}(C)\)，所以 \(-\alpha\in P\cap R_{+}(C')\)，从而 \(\operatorname{Card}(P\cap R_{+}(C'))>\operatorname{Card}(P\cap R_{+}(C))\)。这是荒谬的，因为 \(\alpha\in P\)。因此 \(B(C)\subseteq P\)，进而由命题 19 以及 P 闭这一事实得到 \(R_{+}(C)\subseteq P\)。
\end{enumerate}

推论 1. 设 P 是 R 的子集。以下条件等价：

\begin{enumerate}
\def\labelenumi{(\roman{enumi})}
\item
  存在一个室 C，使得 \(\mathrm{P} = \mathrm{R}_{+}(\mathrm{C})\)；
\item
  P 是闭的，并且 \(\{P, -P\}\) 是 R 的一个划分。
\end{enumerate}

于是，满足 \(\mathrm{P} = \mathrm{R}_{+}(\mathrm{C})\) 的室 C 是唯一的。

若 \(\mathrm{P} = \mathrm{R}_{+}(\mathrm{C})\)，则 \(\mathrm{C}^{-}\) 是 \(x^{*} \in \mathrm{V}^{*}\) 的集合，其中 \(\langle x^{*}, x \rangle > 0\) 对所有 \(x \in \mathrm{P}\) 成立，因而 \(\mathbf{C}\) 是唯一的。

推论 2. 假设 V 具有有序向量空间的结构，使得在此结构下 R 的每个根或者为正或者为负。令 P 为关于此结构的正根集合。则存在唯一的 R 的室 C，使得 \(P = R_{+}(C)\)。

事实上，P 满足推论 1 的条件 (ii)。

这个推论特别适用于所考虑的序是全序的情形；此时关于 R 的条件自动满足。回顾一下，例如选取 V 的一个基 \((e_{i})_{1\leqslant i\leqslant n}\)，并在 V 上取字典序，就可以得到这样的序，于是 \(x=\sum_{i}\xi_{i}e_{i}\) 是 \(\geqslant0\)，如果所有 \(\xi_{i}\) 都为 0，或者 \(\xi_{i}>0\) 对满足 \(\xi_{i}\neq0\) 的最小指标 i 成立。

推论 3. R 的子集 B 是 R 的一个基，当且仅当满足以下条件：

\begin{enumerate}
\def\labelenumi{(\roman{enumi})}
\item
  B 的元素线性无关；
\item
  R 的每个根都是 B 中元素的线性组合，且其系数要么全为正，要么全为负；
\item
  B 的每个根都是不可分的。
\end{enumerate}

我们已经知道这些条件是必要的（第 5 号的 Th. 2 和第 6 号的 Th. 3）。现在假设条件 (i)、(ii)、(iii) 满足。令 P 为由 B 中元素作系数 \(\geqslant0\) 的线性组合而成的根的集合。由于 P 满足推论 1 的条件 (ii)，存在一个室 C，使得 \(P = R_{+}(C)\)；令 \(B' = B(C)\)，令 X 和 \(X'\) 分别为由 B 和 \(B'\) 生成的凸锥。于是

\[
\mathrm{B} \subseteq \mathrm{P} \subseteq \mathrm{X} \quad \text { and } \quad \mathrm{B} ^ {\prime} \subseteq \mathrm{P} \subseteq \mathrm{X} ^ {\prime},
\]

这表明 X 和 \(X'\) 都由 P 生成，因而相等。但是，由 B 中元素（相应地，由 \(B'\) 中元素）生成的半直线是 X（相应地，\(X'\)）的极端生成元；由于这样的半直线只包含一个不可分根，故 \(B = B'\)。

推论 4. 设 B 是 R 的一个基，B' 是 B 的子集，V' 是由 B' 生成的 V 的向量子空间，且 R' = R ∩ V'。则 B' 是根系 R' 的一个基。

这立即由推论 3 和命题 4 的推论得到。

我们称 R' 为由 B' 生成的根系。

推论 5. 设 B 是 R 的一个基，\(A_{1}, A_{2}, \ldots, A_{r}\) 是 B 的两两正交的子集，且 \(A = A_{1} \cup A_{2} \cup \cdots \cup A_{r}\)。则任何由 A 中元素作线性组合得到的根 \(\alpha\)，实际上都是 \(A_{i}\) 中某一个的元素的线性组合。特别地，若 R 不可约，则 B 不存在划分成两两正交子集的划分。

分别令 \(E_1, \ldots, E_r\)、\(E\) 为 \(V\) 中由 \(A_1, \ldots, A_r\)、\(A\) 生成的向量子空间。由推论 4，可假设 \(E = V\)。于是根据第 5 号的 Th. 2 (vii)，各 \(E_i\) 在 \(W(R)\) 下不变，所以 \(R\) 是各 \(R \cap E_i\) 的并集（第 2 号的 Prop. 5）。

推论 6. 采用命题 20 的假设和记号。令 \(V_{1}\) 是由 \(\Sigma\) 生成的 V 的向量子空间。则 \(P \cap (-P) = Q \cup (-Q) = V_{1} \cap R\) 是 \(V_{1}\) 中以 \(\Sigma\) 为基的根系。

我们有 \(P \cap (-P) = (R_{+}(C) \cup Q) \cap ((-R_{+}(C)) \cup (-Q)) = Q \cup (-Q)\)。Th. 3 证明 \(Q \cup (-Q) = V_{1} \cap R\)。最后，根据推论 4，\(\Sigma\) 是根系 \(V_{1} \cap R\) 的一个基。

命题 21. 设 C（相应地 C'）是 R 的一个室，\(\Sigma\)（相应地 \(\Sigma'\)）是 B(C)（相应地 B(C')）的子集，Q（相应地 Q'）是 \(\Sigma\)（相应地 \(\Sigma'\)）中元素以负整数为系数组合成的集合，且 P = Q ∪ R \(_{+}\) (C)（相应地 P' = Q' ∪ R \(_{+}\) (C')）。若存在 Weyl 群中的元素将 P 变为 P'，则存在 Weyl 群中的元素将 C 变为 C' 并将 \(\Sigma\) 变为 \(\Sigma'\)。

我们立即归结到 \(P = P'\) 的情形。令 \(V_1\) 为 \(V\) 中由 \(P \cap (-P)\) 生成的向量子空间。于是 \(\Sigma\) 和 \(\Sigma'\) 是根系 \(R_1 = P \cap (-P)\) 在 \(V_1\) 中的基（命题 20 的推论 6）。因此存在 \(g_1 \in W(R_1)\)，使得 \(g_1(\Sigma) = \Sigma'\)。显然，\(g_1\) 由元素 \(g\) 诱导；该元素属于 \(W(R)\)，且是 \(s_\sigma\)（其中 \(\sigma \in \Sigma\)）的乘积。令 \(\gamma = \sum_{\beta \in B(C)} c_\beta \beta\) 是 \(P - R_1\) 中的元素。于是 \(c_\beta > 0\) 对至少一个 \(\beta \in B(C) - \Sigma\) 成立。此外，若 \(\sigma \in \Sigma\)，则 \(s_\sigma(\gamma) - \gamma \in V_1\)，所以 \(s_\sigma(\gamma)\) 至少有一个系数 \(> 0\)，相对于 \(B(C)\)（第 1 号的公式 (5)），从而 \(s_\sigma(\gamma) \in R_+(C)\)，最终 \(s_\sigma(\gamma) \in P - R_1\)。因此 \(P - R_1\) 在 \(s_\sigma, \sigma \in \Sigma\) 下不变，进而在 \(g\) 下不变，所以 \(g(P) = P\)。于是归结为证明 \(P = P'\) 且 \(\Sigma = \Sigma'\) 的情形。在此情形，\(Q = Q'\)，所以 \(R_+(C) = P - Q = P - Q' = R_+(C')\)，因而 \(C = C'\)（命题 20 的推论 1）。

推论。设 \(\mathrm{P}, \mathrm{P}'\) 是 \(\mathbb{R}\) 的两个抛物子集，它们由 Weyl 群中的元素相互变换。若存在一个室 \(\mathbb{C}\)，属于 \(\mathbb{R}\)，使得 \(\mathrm{R}_{+}(\mathbb{C}) \subseteq \mathrm{P}\) 且 \(\mathrm{R}_{+}(\mathbb{C}) \subseteq \mathrm{P}'\)，则 \(\mathrm{P} = \mathrm{P}'\)。

这由引理 3 和命题 21 得到，因为将 C 变换为 C 的 W(R) 中唯一元素是 1，cf.~第 5 号的 Th. 2。

命题 22. 设 P 是 R 的闭子集，满足 \(P \cap (-P) = \varnothing\)。则存在 R 的一个室 C，使得 \(P \subseteq R_{+}(C)\)。

\begin{enumerate}
\def\labelenumi{\arabic{enumi})}
\item
  根据第 3 号 Th. 1 的 Cor.，假设 \(\alpha \in \mathrm{P}\)、\(\beta \in \mathrm{P}\)、\((\alpha |\beta) < 0\)，则 \(\alpha + \beta \in \mathrm{P}\)。
\item
  我们证明，P 中元素的任何和 \(\alpha_{1} + \cdots + \alpha_{q}\)（\(q \geqslant 1\)）都不为零。我们对 q 归纳。当 q = 1 时断言显然，假设 \(q \geqslant 2\)。若 \(\alpha_{1} + \cdots + \alpha_{q} = 0\)，则
\end{enumerate}

\[
\alpha_ {1} = \alpha_ {2} + \dots + \alpha_ {q},
\]

所以 \((-\alpha_{1}|\alpha_{2}+\cdots+\alpha_{q})>0\)，因而存在 \(j\in[2,q]\)，使得 \((\alpha_{1}|\alpha_{j})<0\)。由证明的第 1 部分，\(\alpha_{1}+\alpha_{j}\in P\)，而关系 \((\alpha_{1}+\alpha_{j})+\sum_{i\neq1,j}\alpha_{i}=0\) 与归纳假设矛盾。

\begin{enumerate}
\def\labelenumi{\arabic{enumi})}
\setcounter{enumi}{2}
\tightlist
\item
  我们证明存在 V 中的非零元素 \(\gamma\)，使得 \((\gamma|\alpha) \geqslant 0\) 对所有 \(\alpha \in P\) 成立。否则，1）的结果将表明，可以找到 P 中元素的一个无限序列 \(\alpha_{1}, \alpha_{2}, \ldots\)，使得
\end{enumerate}

\[
\beta_ {i} = \alpha_ {1} + \dots + \alpha_ {i} \in P
\]

对所有 \(i\) 都有这一点；于是会存在两个不同的整数 \(i, j\)，使得 \(\beta_{i} = \beta_{j}\)，这将与 2）的结果矛盾。

\begin{enumerate}
\def\labelenumi{\arabic{enumi})}
\setcounter{enumi}{3}
\tightlist
\item
  为证明命题，只需（命题 20 的推论 2）证明存在 \((\alpha_{k})_{1\leqslant k\leqslant l}\) 的 \(V\) 的一个基，使得对于由该基定义的字典序，\(P\) 的每个元素都 \(>0\)。我们对 \(l = \dim V\) 归纳，并假设命题已对所有维数 \(< l\) 成立。设 \(\gamma \in V\)，满足 \(\gamma \neq 0\) 且 \((\gamma |\alpha)\geqslant 0\) 对所有 \(\alpha \in P\) 成立（cf.~3））。令 \(L\) 为与 \(\gamma\) 正交的超平面，并令 \(V'\) 为 \(L\) 中由 \(R\cap L\) 生成的子空间。于是 \(R\cap L\) 是 \(V'\) 中的根系，而 \(P\cap L\) 在 \(R\cap L\) 中是闭的。由归纳假设，存在 \((\beta_{1},\ldots ,\beta_{l'})\) 的 \(V'\) 的一个基，使得 \(P\cap L\) 的元素关于该基定义的字典序都是 \(>0\)。于是 \(V\) 的任意基，只要其前 \(l' + 1\) 个元素是 \(\gamma ,\beta_{1},\ldots ,\beta_{l'}\)，其余元素属于 \(L\)，就具有所需性质。
\end{enumerate}

命题 23. 设 P 是 R 的子集，\(V_{1}\)（相应地，\(\Gamma\)）是由 P 生成的 V 的向量子空间（相应地，V 的子群）。以下条件等价：

\begin{enumerate}
\def\labelenumi{(\roman{enumi})}
\item
  P 闭且对称；
\item
  P 闭，且 P 是 \(V_{1}\) 中的根系；
\item
  \(\Gamma \cap \mathbb{R} = \mathbb{P}\)。
\end{enumerate}

假设这些条件满足。对任意 \(\alpha \in \mathbb{P}\)，令 \(\alpha_{1}^{-}\) 为 \(\alpha^{-}\) 在 \(\mathrm{V}_{1}\) 上的限制。于是映射 \(\alpha \mapsto \alpha_{1}^{-}\) 是从根系 \(\mathbb{P}\) 到 \(\mathbb{P}^{-}\) 的典范双射。

\begin{enumerate}
\def\labelenumi{(\roman{enumi})}
\setcounter{enumi}{2}
\item
  \(\Longrightarrow\) (i)：显然。
\item
   \(\Longrightarrow\) (ii)：假设 P 闭且对称。首先，P 满足 \((\mathrm{RS}_{\mathrm{I}})\)，且该根系位于 \(V_{1}\) 中。我们证明，若 \(\alpha, \beta \in P\)，则 \(s_{\alpha}(\beta) \in P\)。当 \(\alpha\) 和 \(\beta\) 成比例时显然成立。否则，\(s_{\alpha}(\beta) = \beta - n(\beta, \alpha)\alpha\)，并且 \(\beta - p\alpha \in R\) 对介于 0 与 \(n(\beta, \alpha)\) 之间的每个有理整数 p 成立（命题 9，第 3 号），所以
\end{enumerate}

\[
\beta - n (\beta , \alpha) \alpha \in P
\]

由于 P 是闭的且对称。因此，\(s_{\alpha,\alpha_{1}}(P) = P\)，且 P 满足 \((RS_{II})\)。显然 P 满足 \((RS_{III})\)。因此，P 满足 (ii)，并且我们同时证明了命题的最后一个断言。

\begin{enumerate}
\def\labelenumi{(\roman{enumi})}
\setcounter{enumi}{1}
\tightlist
\item
  \(\Longrightarrow\) (iii)：我们证明，若条件 (ii) 满足，则 \(\Gamma\cap R=P\)。显然 \(P\subseteq\Gamma\cap R\)。设 \(\beta\in\Gamma\cap R\)。由于 \(\beta\in\Gamma\) 且 \(P=-P\)，有 \(\beta=\alpha_{1}+\alpha_{2}+\cdots+\alpha_{k}\)，其中 \(\alpha_{1},\ldots,\alpha_{k}\in P\)。我们将证明 \(\beta\in P\)。当 k=1 时显然。我们对 k 归纳。我们有
\end{enumerate}

\[
0 <   (\beta | \beta) = \sum_ {i = 1} ^ {k} (\beta | \alpha_ {i}),
\]

所以 \((\beta |\alpha_{i}) > 0\) 对某个指标 \(i\) 成立。若 \(\beta = \alpha_{i}\)，则 \(\beta \in \mathbb{P}\)。否则，\(\beta - \alpha_{i} \in \mathbb{R}\)（第 3 号的 Th. 1 的 Cor.），所以由归纳假设 \(\beta - \alpha_{i} \in \mathbb{P}\)，从而 \(\beta \in \mathbb{P}\)，因为 \(\mathbb{P}\) 是闭的。

Prop. 23 的条件可以在 \(V_{1} = V\) 下实现，但 \(P \neq R\)。例如，当 \(R\) 是 \(G_{2}\) 型的系统而 \(P\) 是 \(A_{2}\) 型的系统时就是这种情况；cf.~Plate X。

命题 24. 设 R' 是 R 与 V 的一个向量子空间的交，从而 R' 是它所生成的向量子空间 V' 中的根系（cf.~命题 4 的推论，第 1 号）。设 B' 是 R' 的一个基。

\begin{enumerate}
\def\labelenumi{(\roman{enumi})}
\item
  存在一个包含 \(\mathbf{B}'\) 的 R 的基。
\item
  \(\mathbf{R}'\) 是 \(\mathbf{R}\) 中由 \(\mathbf{B}'\) 中元素作线性组合得到的元素的集合。
\end{enumerate}

断言 (ii) 显然。我们证明 (i)。设 \((\varepsilon_{1},\varepsilon_{2},\ldots,\varepsilon_{l})\) 是 V 的一个基，使得 \(B' = (\varepsilon_{p+1},\varepsilon_{p+2},\ldots,\varepsilon_{l})\)。与此基对应的 V 上的字典序定义了 R 的一个室 C。显然 \(B'\) 的每个元素在 \(\mathrm{R}_{+}(\mathrm{C})\) 中都是极小的。因此 \(\mathrm{B}' \subseteq \mathrm{B}(\mathrm{C})\)。

